% ID: SEL-D0021-OB-E01
% ESTADO: validada. S022, revisada con los criterios oficiales.
\textbf{a)} El recinto queda por encima de $y=0$ y de
$y=(3-x)/5$, y por debajo de $y=x+3$ y de $y=(30-2x)/7$.
En el eje horizontal resulta $3\le x\le15$. Las fronteras
oblicuas que cortan a $y=x+3$ dan $(-2,1)$ y $(1,4)$.
\begin{center}
\setlength{\unitlength}{4mm}
\begin{picture}(22,10)(-4,-2)
\qbezier(-3,0)(7,0)(17,0)\put(17.2,0){$x$}
\qbezier(0,-1)(0,3)(0,7)\put(0.2,7){$y$}
\qbezier(-2,1)(-0.5,2.5)(1,4)
\qbezier(1,4)(8,2)(15,0)
\qbezier(15,0)(9,0)(3,0)
\qbezier(3,0)(0.5,0.5)(-2,1)
\put(-2,1.4){\makebox(0,0)[r]{$(-2,1)$}}
\put(1,4.4){$(1,4)$}
\put(3,-1){\makebox(0,0){$(3,0)$}}
\put(15,-1){\makebox(0,0){$(15,0)$}}
\end{picture}
\end{center}
\textbf{b)} En $(5,3)$ se cumplen tres inecuaciones, pero
$2\cdot5+7\cdot3=31>30$. Por eso queda fuera.

\textbf{c)} En los cuatro vértices citados, $F=x-y$ vale
$-3,-3,3,15$, respectivamente. El mínimo $\boxed{-3}$ se alcanza
en todo el lado $y=x+3$ con $-2\le x\le1$, pues allí $x-y=-3$.
El máximo $\boxed{15}$ se alcanza sólo en $(15,0)$.
